- w dowodzie twierdzenia o zanurzeniu nerwu są błędy w indeksach. W przedstawieniu punktu należącego do drugiego sympleksu powinno być p_{j_t}, a nie p_{i_t}; trzeba to też poprawić w kolejnych równaniach. Suma odpowiadająca wspólnym wierzchołkom powinna kończyć się na a-1, nie na a. Proszę też dopisać przypadek a=0: wtedy wszystkie współczynniki musiałyby być zerowe, wbrew temu, że ich suma wynosi 1, więc sympleksy są rozłączne.
- we wzorze definiującym funkcję Kuratowskiego trzeba uwzględnić przypadek U_i=X (chodzi o odległość od dopełnienia). Proszę też krótko uzasadnić ciągłość
- w końcowym dowodzie z zawierania się w otwartej kuli promienia \varepsilon/2 wynika średnica <= varepsilon, niekoniecznie < \varepsilon. Najprościej użyć kul promienia \varepsilon/3. Na końcu pokazuje Pani, że funkcja jest \varepsilon-przekształceniem, nie \varepsilon-mapą
- stwierdzenie o homeomorfizmie sympleksu z kostką wymaga uzasadnienia albo odsyłacza -  sama ciągłość rzutowania na współrzędne tego nie pokazuje
- brakuje definicji samej granicy odwrotnej jako podprzestrzeni produktu złożonej z ciągów spełniających \pi_j^i(x_i)=x_j dla j\le i. Proszę też poprawić "zbiór wszystkich dodatnich liczb dodatnich"
- w przykładzie zbioru Cantora proszę krótko opisać homeomorfizm 